% Capítulo de fundamentação: equalização descentralizada sob ruído colorido.
% Comentários CONTEÚDO/TEXTO/ESTADO vêm do esqueleto (rascunho_backbone.tex).

\chapter{Equalização descentralizada sob ruído colorido}\label{chap:dbp-ruido}
% Núcleo do trabalho: o ruído colorido, a estimação da sua covariância, a
% arquitetura descentralizada, o BCD-MMSE e a matriz G_c que cada DU inverte,
% e as técnicas para inverter essa matriz.
% Roteiro T1-T6 de constraints/notas_redacao.md.

\section{\texorpdfstring{\textcolor{red}{Modelo do ruído colorido}}{Modelo do ruído colorido}}
\label{sec:ruido-colorido}
% CONTEÚDO  Do AWGN ao ruído colorido: interferentes + ruído térmico;
%           construção de R = beta Hbar Hbar^H + sigma^2 I; IoT; leitura de R.
% TEXTO     1-noise-and-constraints.tex, "Modelo AWGN" e "Modelo do ruído
%           colorido" (assinatura espacial, hipóteses, construção de R_n).
% ESTADO    escrito até a construção de R; falta a IoT

\textcolor{red}{O equalizador da Equação~{\hypersetup{linkcolor=red}\eqref{eq:lmmse_from_wiener}} depende da matriz de covariância do ruído. Esta seção descreve o ruído considerado neste trabalho e constrói essa matriz.}

\subsection{\texorpdfstring{\textcolor{red}{Modelo AWGN}}{Modelo AWGN}}

\textcolor{red}{No modelo AWGN, em que $\mathbf{R}_n = \sigma^2\mathbf{I}_M$, a covariância do ruído se torna proporcional à matriz identidade. Assim, a solução LMMSE pode ser simplificada definindo o parâmetro de regularização $\rho = \sigma^2/E_s$, que relaciona a potência do ruído com a energia média por símbolo transmitido. Nesse caso, a Equação~\eqref{eq:lmmse_solution} se reduz a~\cite{jeon2019}:}

\begin{equation}
\mathbf{W}_\text{LMMSE}=
\left(
\mathbf{H}^H\mathbf{H}
+
\rho\mathbf{I}
\right)^{-1}
\mathbf{H}^H
\label{eq:lmmse_awgn}
\end{equation}


\textcolor{red}{O equalizador LMMSE tem um pequeno aumento de complexidade computacional em relação ao ZF, pois também exige estimar a matriz de covariância do ruído $\mathbf{R}_n$ para calcular o filtro. Os dois equalizadores têm a mesma complexidade assintótica $\mathcal{O}(K^3)$, dominada pela inversão de matriz. Porém, o LMMSE geralmente oferece melhor desempenho de detecção, pois minimiza a degradação da BER~\cite{li-tradeofs}.}



 

\subsection{\texorpdfstring{\textcolor{red}{Assinatura espacial e interferência}}{Assinatura espacial e interferência}}

O sinal transmitido por um User Equipament (UE) chega em todas as antenas da mesma Estação Base, sendo multiplicado em cada antena pelo respectivo ganho do canal definido em [citar seção do modelo de canal]. O sinal recebido a partir de um único transmissor é escrito como $\mathbf{y}=\mathbf{h}s+\mathbf{n}$, em que o vetor $\mathbf{h}$ agrupa os ganhos de canal entre o UE e cada uma das antenas da Estação Base [ganho precisa ter sido definido antes]. Esse vetor é chamado assinatura espacial do sinal, porque descreve como o sinal se distribui no conjunto de antenas. 

O sinal transmitido por um User Equipament não-alvo da estação base também chega a todas as antenas, com sua própia assinatura de canal em cada uma delas. Dessa forma, a mesma interferência ocorre simultaneamente em várias antenas, com ganhos distintos entre elas. Além dela, cada antena tem o seu próprio ruído térmico, independente entre antenas. \textcolor{red}{A soma dessa interferência com o ruído térmico é chamada, nesse trabalho, de ruído colorido.}

Em cenários com interferência, a relação entre o sinal recebido e o ruído pode ser descrita como:

\begin{equation*}
\mathbf{y}=\mathbf{H s}+\underbrace{\mathbf{n}}_{\text{\textcolor{red}{ruído colorido}}}=\mathbf{H s}+\underbrace{\overline{\mathbf{H}} \overline{\mathbf{s}}}_{\text{\textcolor{red}{interferência}}}+\underbrace{\overline{\mathbf{n}}}_{\text {AWGN }}, \tag{5}
\end{equation*}

Na equação acima, $\overline{\mathbf{H}} \in \mathbb{C}^{M \times r}$ é a matriz de canal entre cada uma das $M$ antenas da estação base e os $r$ UEs (User Equipaments) não-alvo, cada um com uma antena. $\overline{\mathbf{s}} \in \mathcal{S}^{r \times 1}$ representa os símbolos transmitidos pelos $r$ UEs. Já $\overline{\mathbf{n}} \in \mathbb{C}^{M \times 1}$ representa o ruído térmico de fundo de cada antena da estação base, modelado como AWGN de distribuição $\mathcal{CN}(\mathbf{0}, \sigma^2\mathbf{I})$, com média zero e matriz de covariância $\sigma^2\mathbf{I}$.


\begin{figure}[htbp]
\centering
\resizebox{0.7\textwidth}{!}{%
\begin{tikzpicture}[
cell/.style={draw, minimum size=1cm, inner sep=0pt, outer sep=0pt, thick},
brace/.style={decorate, decoration={brace, amplitude=6pt, raise=3pt}, thick}
]

% ==================== Matriz H (M x r) ====================
\begin{scope}[local bounding box=H]
\foreach \row in {1,...,4} {
\foreach \col in {1,...,3} {
\node[cell, fill=cyan!30] (H-\row-\col) at (\col*1, -\row*1) {};
}
}
\end{scope}
\node[above=0.3cm of H] {\Large $\overline{\mathbf{H}} \in \mathbb{C}^{M \times r}$};

% Chaves para H
\draw[brace] ([xshift=-2pt]H.south west) -- ([xshift=-2pt]H.north west) node[midway, xshift=-18pt] {\Large $M$};
\draw[brace] ([yshift=-2pt]H.south east) -- ([yshift=-2pt]H.south west) node[midway, yshift=-18pt] {\Large $r$};

% ==================== Vetor s (r x 1) ====================
\begin{scope}[local bounding box=s, xshift=5cm, yshift=-0.5cm]
\foreach \row in {1,...,3} {
\node[cell, fill=orange!40] (s-\row) at (1, -\row*1) {};
}
\end{scope}
\node[above=0.8cm of s] {\Large $\overline{\mathbf{s}} \in \mathcal{S}^{r \times 1}$};

% Chaves para s
\draw[brace] ([xshift=-2pt]s.south west) -- ([xshift=-2pt]s.north west) node[midway, xshift=-18pt] {\Large $r$};
\draw[brace] ([yshift=-2pt]s.south east) -- ([yshift=-2pt]s.south west) node[midway, yshift=-18pt] {\Large $1$};

% ==================== Vetor n (M x 1) ====================
\begin{scope}[local bounding box=n, xshift=9.5cm]
\foreach \row in {1,...,4} {
\node[cell, fill=yellow!40] (n-\row) at (1, -\row*1) {};
}
\end{scope}
\node[above=0.3cm of n] {\Large $\overline{\mathbf{n}} \in \mathbb{C}^{M \times 1}$};

% Chaves para n
\draw[brace] ([xshift=-2pt]n.south west) -- ([xshift=-2pt]n.north west) node[midway, xshift=-18pt] {\Large $M$};
\draw[brace] ([yshift=-2pt]n.south east) -- ([yshift=-2pt]n.south west) node[midway, yshift=-18pt] {\Large $1$};

\end{tikzpicture}%
}
\caption{Representação das dimensões das matrizes do sistema.}
\label{fig:matrizes}
\end{figure}


\subsection{\texorpdfstring{\textcolor{red}{Hipóteses do modelo}}{Hipóteses do modelo}}

Considera-se que os $r$ UEs não-alvo transmitem de forma independente entre si, com potência $\beta$, de modo que 

$$
\mathbb{E}[\overline{\mathbf{s}}\overline{\mathbf{s}}^H]=\beta\mathbf{I}_r
$$

Assume-se também que o ruído térmico é espacialmente branco, tal que 
$$\mathbb{E}[\overline{\mathbf{n}}\overline{\mathbf{n}}^H]=\sigma^2\mathbf{I}
$$ 

e que $\overline{\mathbf{s}}$ e $\overline{\mathbf{n}}$ são mutuamente independentes e possuem média nula. A matriz $\overline{\mathbf{H}}$ é considerada determinística, uma vez que permanece constante durante dado os mesmos parâmetros e entradas, isto é, no bloco de coerência,
% conforme será explicado posteriormente.
\textcolor{red}{conforme apresentado na Seção~{\hypersetup{linkcolor=red}\ref{subsec:bloco-coerencia}}.}

\begin{figure}[htbp]
\centering
\resizebox{0.8\textwidth}{!}{%
\begin{tikzpicture}[
cell/.style={draw, minimum size=1.2cm, inner sep=0pt, outer sep=0pt, thick},
brace/.style={decorate, decoration={brace, amplitude=6pt, raise=3pt}, thick}
]

% ==================== Matriz de Covariância de s (r x r) ====================
\begin{scope}[local bounding box=covS]
\foreach \row in {1,...,3} {
\foreach \col in {1,...,3} {
\ifnum\row=\col
\node[cell, fill=orange!40] (covS-\row-\col) at (\col*1.2, -\row*1.2) {\Large $\beta$};
\else
\node[cell, fill=black!5] (covS-\row-\col) at (\col*1.2, -\row*1.2) {\large $0$};
\fi
}
}
\end{scope}
\node[above=0.4cm of covS] {\Large $\mathbb{E}[\overline{\mathbf{s}}\overline{\mathbf{s}}^H] = \beta\mathbf{I}_r$};

% Chaves para a matriz de s
\draw[brace] ([xshift=-2pt]covS.south west) -- ([xshift=-2pt]covS.north west) node[midway, xshift=-18pt] {\Large $r$};
\draw[brace] ([yshift=-2pt]covS.south east) -- ([yshift=-2pt]covS.south west) node[midway, yshift=-18pt] {\Large $r$};

% ==================== Matriz de Covariância de n (M x M) ====================
\begin{scope}[local bounding box=covN, xshift=7cm]
\foreach \row in {1,...,4} {
\foreach \col in {1,...,4} {
\ifnum\row=\col
\node[cell, fill=yellow!40] (covN-\row-\col) at (\col*1.2, -\row*1.2) {\Large $\sigma^2$};
\else
\node[cell, fill=black!5] (covN-\row-\col) at (\col*1.2, -\row*1.2) {\large $0$};
\fi
}
}
\end{scope}
\node[above=0.4cm of covN] {\Large $\mathbb{E}[\overline{\mathbf{n}}\overline{\mathbf{n}}^H] = \sigma^2\mathbf{I}$};

% Chaves para a matriz de n
\draw[brace] ([xshift=-2pt]covN.south west) -- ([xshift=-2pt]covN.north west) node[midway, xshift=-18pt] {\Large $M$};
\draw[brace] ([yshift=-2pt]covN.south east) -- ([yshift=-2pt]covN.south west) node[midway, yshift=-18pt] {\Large $M$};

\end{tikzpicture}%
}
\caption{Matrizes de covariância dos sinais não-alvo e do ruído térmico espacialmente branco.}
\label{fig:covariancias}
\end{figure}


\subsection{\texorpdfstring{Construção de $\mathbf{R}_{\textcolor{red}{n}}$}{Construção de Rn}}\label{subsec:construcao-R}
\textcolor{red}{Esta seção constrói, para o cenário com interferência, a matriz $\mathbf{R}_n$ do equalizador LMMSE da Equação~{\hypersetup{linkcolor=red}\eqref{eq:mimo_correlations}}.}

Em uma transformação linear, em que um vetor aleatório $\mathbf{x}$ é obtido a partir de $\mathbf{z}$ por $\mathbf{x} = A\mathbf{z}$, com $A$ determinístico, a matriz de covariância de $\mathbf{x}$ pode ser obtida a partir da matriz de covariância de $\mathbf{z}$ por 

$$
\operatorname{Cov}(\mathbf{x}) = A\operatorname{Cov}(\mathbf{z})A^H,
$$

em que $A^H$ representa a transposta conjugada de $A$. Além disso, considerando duas variáveis aleatórias $\mathbf{u}$ e $\mathbf{v}$ independentes e de média zero, a matriz de covariância do vetor resultante da soma é dada por

$$
\operatorname{Cov}(\mathbf{u}+\mathbf{v})=\operatorname{Cov}(\mathbf{u})+\operatorname{Cov}(\mathbf{v}).
$$

A partir dessas hipóteses e do ruído colorido definido como $\mathbf{n}=\overline{\mathbf{H}}\overline{\mathbf{s}}+\overline{\mathbf{n}}$, as matrizes de covariância das parcelas $\overline{\mathbf{H}}\overline{\mathbf{s}}$ e $\overline{\mathbf{n}}$ podem ser somadas, pois são independentes e possuem média nula. Dessa forma,

$$
\mathbf{R}_{\textcolor{red}{n}}\triangleq\mathbb{E}[\mathbf{n}\mathbf{n}^H]
=\operatorname{Cov}(\overline{\mathbf{H}}\overline{\mathbf{s}})
+\operatorname{Cov}(\overline{\mathbf{n}}).
$$

A covariância da primeira parcela pode ser obtida pela propriedade da transformação linear, considerando $A=\overline{\mathbf{H}}$ e $\mathbf{z}=\overline{\mathbf{s}}$. Assim,

$$
\operatorname{Cov}(\overline{\mathbf{H}}\overline{\mathbf{s}})
=\overline{\mathbf{H}}\mathbb{E}[\overline{\mathbf{s}}\overline{\mathbf{s}}^H]\overline{\mathbf{H}}^H
=\beta\overline{\mathbf{H}}\overline{\mathbf{H}}^H.
$$

Como $\operatorname{Cov}(\overline{\mathbf{n}})=\sigma^2\mathbf{I}$, obtém-se, portanto, a matriz de covariância do ruído colorido:

$$
\mathbf{R}_{\textcolor{red}{n}}=\beta\overline{\mathbf{H}}\overline{\mathbf{H}}^H+\sigma^2\mathbf{I}.
$$

\textcolor{red}{Assim, no cenário com interferência, a matriz $\mathbf{R}_{yy}$ da Equação~{\hypersetup{linkcolor=red}\eqref{eq:mimo_correlations}} fica}

$$
\textcolor{red}{\mathbf{R}_{yy}=E_s\mathbf{H}\mathbf{H}^H+\beta\overline{\mathbf{H}}\overline{\mathbf{H}}^H+\sigma^2\mathbf{I}.}
$$


\begin{figure}[htbp]
\centering
\resizebox{\textwidth}{!}{%
\begin{tikzpicture}[
    cell/.style={draw, minimum size=1cm, inner sep=0pt, outer sep=0pt, thick},
    brace/.style={decorate, decoration={brace, amplitude=6pt, raise=3pt}, thick}
]

% ==================== Parêntese Esquerdo ====================
\node[scale=6] at (0, -2.5) {$($};

% ==================== 1. Matriz H (M x r) ====================
\begin{scope}[local bounding box=Hbox, xshift=1cm]
    \foreach \row in {1,...,4} {
        \foreach \col in {1,...,3} {
            \node[cell, fill=cyan!30] (H-\row-\col) at (\col*1, -\row*1) {};
        }
    }
\end{scope}
\node[above=0.4cm of Hbox] {\Large $\overline{\mathbf{H}}$};
\draw[brace] ([xshift=-2pt]Hbox.south west) -- ([xshift=-2pt]Hbox.north west) node[midway, xshift=-18pt] {\Large $M$};
\draw[brace] ([yshift=-2pt]Hbox.south east) -- ([yshift=-2pt]Hbox.south west) node[midway, yshift=-18pt] {\Large $r$};

% ==================== 2. Matriz Beta (r x r) ====================
% xshift ajustado para 5.5cm (garante espaço perfeito para a chave esquerda)
\begin{scope}[local bounding box=betaBox, xshift=5.5cm, yshift=-0.5cm]
    \foreach \row in {1,...,3} {
        \foreach \col in {1,...,3} {
            \ifnum\row=\col
                \node[cell, fill=orange!40] (B-\row-\col) at (\col*1, -\row*1) {\Large $\beta$};
            \else
                \node[cell, fill=black!5] (B-\row-\col) at (\col*1, -\row*1) {\large $0$};
            \fi
        }
    }
\end{scope}
\node[above=0.4cm of betaBox] {\Large $\beta\mathbf{I}_r$};
\draw[brace] ([xshift=-2pt]betaBox.south west) -- ([xshift=-2pt]betaBox.north west) node[midway, xshift=-18pt] {\Large $r$};
\draw[brace] ([yshift=-2pt]betaBox.south east) -- ([yshift=-2pt]betaBox.south west) node[midway, yshift=-18pt] {\Large $r$};

% ==================== 3. Matriz H^H (r x M) ====================
% xshift ajustado para 10cm
\begin{scope}[local bounding box=HHbox, xshift=10cm, yshift=-0.5cm]
    \foreach \row in {1,...,3} {
        \foreach \col in {1,...,4} {
            \node[cell, fill=cyan!30] (HH-\row-\col) at (\col*1, -\row*1) {};
        }
    }
\end{scope}
\node[above=0.4cm of HHbox] {\Large $\overline{\mathbf{H}}^H$};
\draw[brace] ([xshift=-2pt]HHbox.south west) -- ([xshift=-2pt]HHbox.north west) node[midway, xshift=-18pt] {\Large $r$};
\draw[brace] ([yshift=-2pt]HHbox.south east) -- ([yshift=-2pt]HHbox.south west) node[midway, yshift=-18pt] {\Large $M$};

% ==================== Parêntese Direito ====================
\node[scale=6] at (15.5, -2.5) {$)$};

% ==================== Operador de Soma (+) ====================
\node at (17, -2.5) {\Huge $+$};

% ==================== 4. Matriz de Ruído Térmico (M x M) ====================
% xshift ajustado para 18cm
\begin{scope}[local bounding box=covN, xshift=18cm]
    \foreach \row in {1,...,4} {
        \foreach \col in {1,...,4} {
            \ifnum\row=\col
                \node[cell, fill=yellow!40] (t2-\row-\col) at (\col*1, -\row*1) {\Large $\sigma^2$};
            \else
                \node[cell, fill=black!5] (t2-\row-\col) at (\col*1, -\row*1) {\large $0$};
            \fi
        }
    }
\end{scope}
\node[above=0.4cm of covN] {\Large $\sigma^2\mathbf{I}$};
\draw[brace] ([xshift=-2pt]covN.south west) -- ([xshift=-2pt]covN.north west) node[midway, xshift=-18pt] {\Large $M$};
\draw[brace] ([yshift=-2pt]covN.south east) -- ([yshift=-2pt]covN.south west) node[midway, yshift=-18pt] {\Large $M$};

% ==================== Operador de Igualdade (=) ====================
\node at (23.5, -2.5) {\Huge $=$};

% ==================== 5. Matriz R (M x M) ====================
% xshift ajustado para 24.5cm
\begin{scope}[local bounding box=resultR, xshift=24.5cm]
    \foreach \row in {1,...,4} {
        \foreach \col in {1,...,4} {
            \node[cell, fill=green!30] (R-\row-\col) at (\col*1, -\row*1) {};
        }
    }
\end{scope}
\node[above=0.4cm of resultR] {\Large $\mathbf{R}_{\textcolor{red}{n}}$};
\draw[brace] ([xshift=-2pt]resultR.south west) -- ([xshift=-2pt]resultR.north west) node[midway, xshift=-18pt] {\Large $M$};
\draw[brace] ([yshift=-2pt]resultR.south east) -- ([yshift=-2pt]resultR.south west) node[midway, yshift=-18pt] {\Large $M$};

\end{tikzpicture}%
}
\caption{Formação da matriz de covariância do ruído colorido $\mathbf{R}_{\textcolor{red}{n}}$ a partir de seus fatores aditivos e multiplicativos.}
\label{fig:matriz_R}
\end{figure}

\subsection{Interferência sobre ruído térmico (IoT)}
% Novo: definição da IoT; papel de beta.

\subsection{\texorpdfstring{\textcolor{red}{O que $\mathbf{R}_n$ representa}}{O que Rn representa}}\label{subsec:o-que-R-representa}

% [ressaltar que a matriz de ruido eh MxM]

%% assim como tem a imagem das matrizes de covariancia e do ruido termico, mostrar como fica a matriz final, possivelmente fazer uma soma na figura dois 
A matriz $\mathbf{R}_{\textcolor{red}{n}}$ é composta por duas contribuições com características distintas. O termo $\sigma^2\mathbf{I}$ representa o ruído térmico. Como sua matriz de covariância possui apenas elementos na diagonal, o ruído presente em diferentes cadeias de recepção não é correlacionado. Já o termo $\beta\overline{\mathbf{H}}\overline{\mathbf{H}}^H$ representa a contribuição dos sinais transmitidos pelos UEs não-alvo após sua propagação pelo canal. Como cada um desses sinais é recebido por múltiplas antenas, o mesmo sinal pode contribuir para diferentes componentes do vetor recebido, estabelecendo uma relação entre elas. Dessa forma, essa parcela pode apresentar elementos fora da diagonal e introduzir correlação entre as antenas.

\section{\texorpdfstring{\textcolor{red}{Estimação de $\mathbf{R}_n$}}{Estimação de Rn}}
\label{sec:estimacao-R}
% CONTEÚDO  Roteiro T1: R não pode ser calculada; estimativa amostral com N
%           amostras de piloto; posto de R-chapéu; quantidade e custo de N.
% TEXTO     1-noise-and-constraints.tex, "Estimação de R_n" (subseções
%           parciais); constraints/rascunho_filtros_covariancia.tex.
% ESTADO    parcial

\subsection{\texorpdfstring{\textcolor{red}{Por que estimar $\mathbf{R}_n$}}{Por que estimar Rn}}

A matriz $\mathbf{R}_{\textcolor{red}{n}}$ descreve a matriz de covariância do ruído colorido completo, ou seja, a soma da interferência com o ruído térmico. Entretanto, não é possível obtê-la praticamente:  $\mathbb{E}[\cdot]$ é definida como uma média sobre infinitas realizações, desse modo, não é possível avaliar \textcolor{red}{$\mathbb{E}[\mathbf{n}\mathbf{n}^H]$} com um número finito de observações; além disso, $
\mathbf{R}_{\textcolor{red}{n}}=\beta\overline{\mathbf{H}}\overline{\mathbf{H}}^H+\sigma^2\mathbf{I}
$ depende das matrizes de canal $\overline{\mathbf{H}}$ das UEs não alvo, que não é estimada pelo receptor, assim como a potência $\beta$ de transmissão dos usuários não é. Dessa forma, $\mathbf{R}_{\textcolor{red}{n}}$ precisa ser estimada a partir de observações do próprio ruído. Essa estimativa é obtida através da média amostral

$$\hat{\mathbf{R}}_{\textcolor{red}{n}}=\frac{1}{N}\sum_{i=1}^{N}\mathbf{n}^i(\mathbf{n}^i)^H$$

em que \sout{$\mathbf{N}$} \textcolor{red}{$N$} é o número de amostras de ruído disponíveis.

\subsection{\texorpdfstring{\textcolor{red}{Como as amostras são obtidas}}{Como as amostras são obtidas}}
% Absorve a seção "Pilotos" (vazia) do capítulo atual.

As amostras de ruído são obtidas durante a etapa de símbolos piloto, sendo $\mathbf{n}^i \in \mathbb{C}^{M\times1}$ a amostra de ruído observada no $i$-ésimo elemento de recurso (RE) de piloto.  Nessa etapa, o sinal transmitido pelos UEs alvo é conhecido pelo receptor, e sua contribuição pode ser cancelada do sinal recebido. O que resta é o valor da interferência somada ao ruído térmico. 

\textcolor{red}{Cada amostra $\mathbf{n}^i$ é um vetor com um valor de ruído colorido para cada antena \cite{zhao2025jsac}. Assim, cada amostra é uma realização do mesmo ruído descrito pela matriz $\mathbf{R}_n$.}
% FONTE: artigo base, paper.md:123 (n^i em C^{M x 1}, amostra do i-ésimo RE de piloto).

\subsection{\texorpdfstring{\textcolor{red}{Intervalo das amostras}}{Intervalo das amostras}}

\textcolor{red}{Na frequência, a covariância do ruído permanece constante ao longo de um intervalo largo, e por isso pode ser estimada com uma média ao longo da frequência \cite{zhao2025jsac}. Assim, as amostras podem ser colhidas em REs de piloto de subportadoras diferentes. Para o canal, as respostas em frequências afastadas deixam de ser correlacionadas, e medir o canal em um conjunto suficiente de frequências permite estimar a sua covariância \cite{barriac-madhow}. Porém, é preciso haver subportadoras suficientes para que a estimativa se aproxime da covariância real \cite{barriac-madhow}.}
% FONTE: artigo base, nota 4 (paper.md:821, chamada em :115), que cita [29].
%        [29] Barriac e Madhow, pilot-estimation/Space*.md:96 e :109 (resultado para a
%        covariância do CANAL; a transferência para o ruído é feita pelo artigo base).

A média amostral é realizada uma vez por bloco de coerência e reutilizada ao longo dos $N_{\text{coh}}$ símbolos subsequentes, em que se assume que a covariância permanece estática.  

\textcolor{red}{Nesse intervalo, o canal e a covariância do ruído são considerados estáticos, e ambos mudam a cada novo bloco de coerência \cite{zhao2025jsac}. A estimativa obtida com os pilotos é usada para processar os sinais de dados recebidos em seguida \cite{tomasin-interference}. Para isso, a covariância precisa variar devagar o suficiente para que a estimativa ainda seja válida quando for usada \cite{barriac-madhow}.}
% FONTE: artigo base, paper.md:561; Tomasin et al. [22], pilot-estimation/Estimation*.md:68;
%        Barriac e Madhow [29], pilot-estimation/Space*.md:109, item (i).

\textcolor{red}{Dessa forma, as $N$ amostras são colhidas em REs de piloto dentro de um mesmo bloco de coerência. A estimativa $\hat{\mathbf{R}}_n$ obtida com elas vale para os $N_{\text{coh}}$ símbolos desse bloco.}
% FONTE: síntese de paper.md:123 (REs de piloto) e :561 (tempo).

\subsection{\texorpdfstring{\textcolor{red}{Uma estimativa por PRB}}{Uma estimativa por PRB}}\label{subsec:estimativa-prb}

\textcolor{red}{No artigo base, o PRB é a faixa de frequência e tempo em que o canal fica constante \cite{zhao2025jsac}. A matriz $\mathbf{R}_n = \beta\overline{\mathbf{H}}\overline{\mathbf{H}}^H+\sigma^2\mathbf{I}$ depende do canal $\overline{\mathbf{H}}$ dos UEs não-alvo \cite{zhao2025jsac}. Neste trabalho, considera-se que o canal dos UEs não-alvo também fica constante dentro do PRB. Assim, $\mathbf{R}_n$ também fica constante dentro do PRB, e muda de um PRB para outro.}
% FONTE: paper.md:614 (definição de PRB); paper.md:113 (R depende de Hbar).
% HIPÓTESE nossa: os interferentes têm o mesmo Doppler da condição de teste.

\textcolor{red}{Por isso, este trabalho usa uma estimativa $\hat{\mathbf{R}}_n$ para cada PRB, feita com $N$ amostras do próprio PRB. Esse é o modelo das simulações do artigo base, em que cada realização do canal tem $N = 192$ amostras de ruído \cite{zhao2025jsac}. O simulador deste trabalho segue o mesmo modelo: a cada realização, ele gera um novo $\overline{\mathbf{H}}$ e $N$ novas amostras de ruído.}
% FONTE: paper.md:577 (N = 192 por realização); simulador: dbp_sim/m1_driver.py:163-178
%        (Hbar e noise_samples gerados dentro do laço de realizações).

\textcolor{red}{A média ao longo da frequência, citada acima, seria uma alternativa: uma única estimativa para a banda inteira. Ela mudaria as curvas de SER, e por isso não é usada nos resultados.}
% FONTE: nota 4, paper.md:821. Decisão da autora (2026-10-12): modelo B (por PRB).
%        O modelo A fica como teste extra opcional (software/simulador_modelo_a).

\textcolor{red}{No tempo, o bloco de coerência tem a duração do tempo de coerência $T_c$ do canal. Neste trabalho, $T_c$ é calculado a partir do Doppler máximo $f_D$ de cada condição de teste \cite{3gpp-ts38104}, por $T_c = 1/(8 f_D)$ \cite[Eq.~2.44]{tse-viswanath}. As matrizes $\hat{\mathbf{R}}_n$, $\mathbf{G}_c$ e $\mathbf{W}$ são refeitas uma vez a cada $T_c$. Assim, $N_{\text{coh}}$ é o número de símbolos OFDM em $T_c$.}
% FONTE: f_D: [R5] TS 38.104 v16.4.0, Anexo G, Tab. G.2.2-1, G.2.2-2 (p. 241), G.3-1 e
%        G.3-2 (p. 252); T_c = 1/(4 D_s) = 1/(8 f_D): Tse e Viswanath, Eq. (2.44), p. 31.
%        Refazer a cada T_c: decisão da autora, opção (b), apoiada em paper.md:614.
%        Valores de T_c e símbolos por T_c: software/Restrições - Parâmetros.csv, linhas 17-24.

\subsection{\texorpdfstring{\textcolor{red}{O que cada amostra contribui}}{O que cada amostra contribui}}

\textcolor{red}{A Figura~{\hypersetup{linkcolor=red}\ref{fig:posto_unitario}} mostra a parcela $\mathbf{n}^i(\mathbf{n}^i)^H$ referente a uma única amostra de ruído.}

\begin{figure}[htbp]
\centering
\resizebox{\textwidth}{!}{%
\begin{tikzpicture}[
    cell/.style={draw, minimum size=1.2cm, inner sep=0pt, outer sep=0pt, thick},
    brace/.style={decorate, decoration={brace, amplitude=6pt, raise=3pt}, thick}
]

% ==================== 1. Vetor Coluna n^i (M x 1) ====================
\begin{scope}[local bounding box=ni]
    \foreach \row in {1,...,4} {
        \node[cell, fill=yellow!40] (ni-\row) at (0, -\row*1.2) {\large $n_\row$};
    }
\end{scope}
\node[above=0.4cm of ni] {\Large $\mathbf{n}^i$};

\draw[brace] ([xshift=-2pt]ni.south west) -- ([xshift=-2pt]ni.north west) node[midway, xshift=-18pt] {\Large $M$};
\draw[brace] ([yshift=-2pt]ni.south east) -- ([yshift=-2pt]ni.south west) node[midway, yshift=-18pt] {\Large $1$};

% ==================== 2. Vetor Linha (n^i)^H (1 x M) ====================
\begin{scope}[local bounding box=nih, xshift=2cm]
    \foreach \col in {1,...,4} {
        % O vetor linha é centralizado verticalmente em relacao ao vetor coluna (y=-3.0 corresponde ao meio)
        \node[cell, fill=yellow!40] (nih-\col) at (\col*1.2, -3.0) {\large $n_\col^*$}; 
    }
\end{scope}
\node[above=0.4cm of nih] {\Large $(\mathbf{n}^i)^H$};

\draw[brace] ([xshift=-2pt]nih.south west) -- ([xshift=-2pt]nih.north west) node[midway, xshift=-18pt] {\Large $1$};
\draw[brace] ([yshift=-2pt]nih.south east) -- ([yshift=-2pt]nih.south west) node[midway, yshift=-18pt] {\Large $M$};

% ==================== Operador de Igualdade (=) ====================
\node at (8.0, -3.0) {\Huge $=$};

% ==================== 3. Matriz Resultante (M x M) ====================
\begin{scope}[local bounding box=ni_result, xshift=8.5cm]
    \foreach \row in {1,...,4} {
        \foreach \col in {1,...,4} {
            \node[cell, fill=yellow!40] (res-\row-\col) at (\col*1.2, -\row*1.2) {\normalsize $n_\row n_\col^*$};
        }
    }
\end{scope}
\node[above=0.4cm of ni_result] {\Large $\mathbf{n}^i(\mathbf{n}^i)^H$};

\draw[brace] ([xshift=-2pt]ni_result.south west) -- ([xshift=-2pt]ni_result.north west) node[midway, xshift=-18pt] {\Large $M$};
\draw[brace] ([yshift=-2pt]ni_result.south east) -- ([yshift=-2pt]ni_result.south west) node[midway, yshift=-18pt] {\Large $M$};

\end{tikzpicture}%
}
\caption{Dimensões do produto externo para a parcela da estimativa da matriz de covariância.}
\label{fig:posto_unitario}
\end{figure}

Como cada celula representa
% a interfencia de uma antena em outra,
\textcolor{red}{o produto $n_a n_b^*$ entre o ruído observado na antena $a$ e o ruído observado na antena $b$ em uma mesma amostra, a média dessas células ao longo das $N$ amostras se aproxima do elemento correspondente de $\mathbf{R}_n$, que indica o quanto o ruído dessas duas antenas varia em conjunto.}

O posto de uma matriz é o número máximo de linhas ou colunas linearmente independentes que essa matriz possui.  Cada parcela $\mathbf{n}^i(\mathbf{n}^i)^H$ é uma matriz de posto unitário, uma vez que todas as suas colunas são múltiplas do mesmo vetor $\mathbf{n}^i$.

\textcolor{red}{Uma única amostra mostra apenas uma realização do ruído. Calcular o produto para várias realizações independentes e tirar a média desses produtos resulta em uma estimativa da correlação \cite[p.~229]{haykin-sistemas}. Essa estimativa se aproxima de $\mathbf{R}_n$ conforme o número de amostras cresce, mas essa convergência é lenta \cite{tomasin-interference}.}
% FONTE: livro de probabilidade, Exemplo 5.16, p. 229; Tomasin et al. [22], Estimation*.md:109.

\subsection{\texorpdfstring{\textcolor{red}{Quantidade de amostras}}{Quantidade de amostras}}
% Posto de R-chapéu <= N: motiva o condicionamento (subseção seguinte).

\textcolor{red}{O erro quadrático médio de cada elemento de $\hat{\mathbf{R}}_n$ diminui na proporção de $1/N$, considerando amostras gaussianas e independentes entre si \cite{tomasin-interference}. Assim, quanto maior o número de amostras, mais precisa é a estimativa. Por isso, o artigo base exige $N \gg K$, para garantir a acurácia da estimativa \cite{zhao2025jsac}. Quando $N$ é grande, o erro entre $\hat{\mathbf{R}}_n$ e $\mathbf{R}_n$ é pequeno \cite{zhao2025jsac}.}
% FONTE: Tomasin et al. [22], eq. (21)-(23), Estimation*.md:242-259 (hipóteses em :57 e :246);
%        artigo base, paper.md:123 (N >> K) e nota 5, paper.md:821.

\textcolor{red}{Nas simulações, o artigo base usa $N = 192$ \cite{zhao2025jsac}. Um trabalho anterior dos mesmos autores cita 96 ou 192 como valores típicos de $N$ \cite{zhao2021-globecom}. Como as amostras vêm dos REs de piloto do PRB, $N$ também fica limitado ao número de REs de piloto que cabem em $T_c$.}
% FONTE: paper.md:577 e :640 (N = 192); zhao2021_globecom, linha 54 ("typically 96 or 192").
% NOTA: a última frase é dedução (paper.md:123 + :614).

\subsection{\texorpdfstring{\textcolor{red}{REs de piloto disponíveis}}{REs de piloto disponíveis}}\label{subsec:res-piloto}

\textcolor{red}{Neste trabalho, as amostras de ruído vêm dos REs do sinal de referência de demodulação (DM-RS). Um PRB tem 12 subportadoras \cite[Sec.~4.4.4.1]{3gpp-ts38211}, e um slot tem 14 símbolos OFDM \cite[Tab.~4.3.2-1]{3gpp-ts38211}. O DM-RS ocupa no máximo 4 símbolos por slot \cite[Tab.~6.4.1.1.3-3 e 6.4.1.1.3-4]{3gpp-ts38211}. Assim, cabem no máximo $4 \times 12 = 48$ REs de DM-RS por PRB em um slot. Para juntar $N = 192$ amostras, são precisos pelo menos 4 slots.}
% FONTE: [R2] TS 38.211 v18.6.0: Sec. 4.4.4.1, p. 15 (12 subportadoras); Tab. 4.3.2-1,
%        p. 14 (14 símbolos por slot); Tab. 6.4.1.1.3-3 e -4, p. 101 (até 4 símbolos de DM-RS).
% HIPÓTESE nossa: contar só o DM-RS. O artigo diz apenas "pilot REs" (paper.md:123).
%        O SRS poderia dar mais REs ([R2] Sec. 6.4.1.4.1, p. 109-110), mas não é usado.

\textcolor{red}{A Tabela~{\hypersetup{linkcolor=red}\ref{tab:n-maximo}} mostra quantos REs de DM-RS cabem em um PRB dentro de $T_c$, para cada condição de teste. Nas condições mais lentas, as 192 amostras cabem em $T_c$. Nas mais rápidas, não cabem, e $N$ fica limitado ao número máximo de REs. Nos trens de alta velocidade, cabe só um símbolo OFDM inteiro em $T_c$, e restam 12 REs.}
% FONTE: software/Restrições - Parâmetros.csv, tabela "Número máximo de REs de piloto
%        (DM-RS) por PRB em T_c" (linhas 44-53). Símbolos de DM-RS em S símbolos:
%        4*floor(S/14) + min(S mod 14, 4).

\begin{table}[htbp]
\caption{\textcolor{red}{Número máximo de REs de DM-RS por PRB dentro de $T_c$.}}
\label{tab:n-maximo}
\centering
\footnotesize
\renewcommand{\arraystretch}{1.2}
\begin{tabular}{|l|r|r|r|r|}
\hline
Condição de teste & $T_c$ (ms) & Símbolos OFDM em $T_c$ & REs de DM-RS (máx.) & $N$ usado \\
\hline
TDLA30-5 (FR1)    & 25,0000 & 700 & 2400 & 192 \\
TDLA30-10 (FR1)   & 12,5000 & 350 & 1200 & 192 \\
TDLA30-75 (FR2)   & 1,6667  & 186 & 672  & 192 \\
TDLA30-300 (FR2)  & 0,4167  & 46  & 192  & 192 \\
TDLC300-100 (FR1) & 1,2500  & 35  & 144  & 144 \\
TDLB100-400 (FR1) & 0,3125  & 8   & 48   & 48 \\
HST350 (FR1)      & 0,0536  & 1   & 12   & 12 \\
HST500 (FR1)      & 0,0375  & 1   & 12   & 12 \\
\hline
\end{tabular}
\end{table}
% FONTE da tabela: f_D e condições: [R5] Anexo G; T_c: Tse e Viswanath, Eq. (2.44);
%        símbolo OFDM: 1/28 ms (30 kHz) e 1/112 ms (120 kHz), [R2] Tab. 4.3.2-1, p. 14.

\textcolor{red}{Essa limitação vale para um receptor em tempo real. A simulação não tem tempo: cada realização é um instante isolado, com $N$ amostras próprias. Por isso, as curvas de SER com $N = 192$ continuam válidas como no artigo base. Para ver o efeito de um $N$ menor, este trabalho propõe testes extras com $N = 12$, 48, 96, 144 e 192.}
% FONTE: dbp_sim/m1_driver.py:163-178 (sem eixo de tempo); dbp_sim/params.py:302 (N = 192).
% ESTADO: testes com N diferente ainda não rodados (falta a opção --n-noise-samples).

\subsection{Condicionamento da estimativa}\label{subsec:condicionamento-R}
% CONTEÚDO  Quando N < M, o posto de R-chapéu é no máximo N e a matriz não
%           tem inversa (M = 256, N = 192). O LMMSE precisa dessa inversa.
%           Shrinkage de Ledoit-Wolf: R-chapéu vira (1 - rho) R-chapéu +
%           rho (tr R-chapéu / M) I, com rho calculado a partir das próprias
%           amostras (sem constante a escolher). Alternativa: carga diagonal
%           fixa, que depende de um epsilon.
% TEXTO     nenhum na fundamentação; parágrafo sobre Ledoit-Wolf em
%           rascunho_setup_simulacao.tex (subseção Receptor);
%           dbp_sim/params.py, M6_Covariance.shrinkage.
% FONTE     Ledoit e Wolf (2004), J. Multivariate Analysis 88(2) -- falta
%           no .bib; posto <= N: livro-texto de álgebra linear.
% MÓDULO    M6 (Conditioner em dbp_sim/m5_noise.py)
% ESTADO    a escrever

\textcolor{red}{Cada parcela $\mathbf{n}^i(\mathbf{n}^i)^H$ tem posto unitário. Por isso, a soma de $N$ parcelas tem posto de no máximo $N$. Nas DUs, a estimativa de cada bloco tem tamanho $M_c \times M_c$ (Seção~{\hypersetup{linkcolor=red}\ref{subsec:estimacao-dus}}). Com $N = 192$, o posto não limita o bloco, pois $M_c$ vale 8, 16 ou 32 \cite{zhao2025jsac}. Já nas condições rápidas da Tabela~{\hypersetup{linkcolor=red}\ref{tab:n-maximo}}, $N$ pode ser 12. Nesse caso, um bloco com $M_c = 16$ tem posto de no máximo 12 e não tem inversa. Além disso, $N = 12$ não atende $N \gg K$, com $K = 8$ \cite{zhao2025jsac}. Este trabalho usa o encolhimento de Ledoit-Wolf em cada bloco, que mantém a matriz inversível \cite{ledoit-wolf}.}
% FONTE: posto unitário: parágrafo anterior desta seção; M_c = M/C com M = 128 ou 256 e
%        C = 8 ou 16: paper.md:577; N >> K: paper.md:123; Ledoit-Wolf por bloco:
%        dbp_sim/m5_noise.py:159-168 (Conditioner.covariance), dbp_sim/m7_equalizers.py:71,
%        dbp_sim/params.py:334.
% NOTA: quando N < M_c, o simulador não calcula rho no bloco; usa a intensidade calculada
%        com a matriz de amostras inteira (m5_noise.py:154-157 e :166-167).
% FALTA no biblio.bib: ledoit-wolf (Ledoit e Wolf, 2004, J. Multivariate Analysis 88(2)).

\subsection{\texorpdfstring{\textcolor{red}{Custo do N}}{Custo do N}}

\textcolor{red}{Por outro lado, um $N$ maior aumenta a quantidade de dados que precisa ser guardada e transmitida. Cada DU guarda as $N$ amostras de ruído das suas antenas \cite{zhao2025jsac}. Trocar essas amostras entre as DUs exigiria transferir dados de tamanho $M \times N$ \cite{zhao2025jsac}. No BCD-MMSE, a banda de comunicação em cada iteração ainda depende de $N$ \cite{zhao2025jsac}. Como a estimativa só muda a cada bloco de coerência, esse custo é pago uma vez a cada $N_{\text{coh}}$ símbolos \cite{zhao2025jsac}.}
% FONTE: artigo base, paper.md:493 (N_c em C^{M_c x N} guardado na DU c); :130 (M x N);
%        :471 (banda por iteração depende de N); :561 e :571 (termos divididos por N_coh).

\textcolor{red}{Um $N$ maior também leva mais tempo para ser juntado. Com 48 REs de DM-RS por PRB em cada slot (Seção~{\hypersetup{linkcolor=red}\ref{subsec:res-piloto}}), o tempo de coleta das $N$ amostras é}
\begin{equation*}
\textcolor{red}{t_{\text{coleta}} = \min\left(T_c;\ \left\lceil \frac{N}{48} \right\rceil \text{ slots}\right).}
\end{equation*}
\textcolor{red}{Com $N = 192$, são 4 slots: 2~ms em FR1, com slots de 0,5~ms, e 0,5~ms em FR2, com slots de 0,125~ms \cite[Tab.~4.3.2-1]{3gpp-ts38211}. Nas condições rápidas, as amostras ocupam o $T_c$ inteiro.}
% FONTE: slot de 0,5 ms (mu = 1, 30 kHz) e 0,125 ms (mu = 3, 120 kHz): [R2] Tab. 4.3.2-1,
%        p. 14. Regra de t_coleta: conta nossa; planilha, linhas 67-76.

\textcolor{red}{A estimativa só fica pronta depois da coleta. Por isso, este trabalho considera um processamento em etapas (\textit{pipeline}). Enquanto junta as amostras de um bloco, a DU guarda os dados recebidos nesse bloco. Depois, calcula $\hat{\mathbf{R}}_n$, $\mathbf{G}_c$ e $\mathbf{W}$ e equaliza os dados guardados. Ao mesmo tempo, já junta as amostras do bloco seguinte. Assim, o cálculo de cada bloco precisa terminar em até $T_c$. O tempo total, da primeira amostra até o resultado, é $t_{\text{coleta}} + t_{\text{proc}}$, em que $t_{\text{proc}}$ é o tempo de cálculo.}
% FONTE: decisão da autora (2026-10-12). Condições t_proc <= T_c e t_coleta + t_proc <= D:
%        planilha, tabela "Fórmulas de latência". D fica só como fórmula (sem fonte).
%        Os dados guardados seguem paper.md:561 (W reutilizado nos N_coh símbolos).

\section{Definição Decentralizada}
\label{sec:dbp}
% CONTEÚDO  Arquitetura DBP: clusters, DUs, CU; estrela e daisy chain;
%           roteiro T2 (o que cada DU alcança) e T3 (a bifurcação: descartar
%           os blocos cruzados -> BDAC; mantê-los comprimidos -> BCD).
% TEXTO     1-noise-and-constraints.tex, "Definição Decentralizada" (vazia)
%           e "Estimação nas DUs" (2 linhas).
% ESTADO    vazio



\subsection{Arquitetura e particionamento das antenas}

\subsection{\texorpdfstring{\textcolor{red}{Estimação nas DUs}}{Estimação nas DUs}}\label{subsec:estimacao-dus}

\textcolor{red}{Conforme apresentado na Seção~{\hypersetup{linkcolor=red}\ref{subsec:o-que-R-representa}}, a parcela $\beta\overline{\mathbf{H}}\overline{\mathbf{H}}^H$ introduz correlação entre as antenas.}
%% TIRAR DAQUI E DILUIR NO RESTO DO TESTO OU POR DEPOIS 
Essa correlação tem uma consequência importante para a arquitetura considerada. Ao particionar as antenas em diferentes clusters associados às DUs, a matriz $\mathbf{R}_{\textcolor{red}{n}}$, em geral, não possui estrutura bloco-diagonal. Assim, os blocos fora da diagonal, dados por $\mathbf{R}_{mn}=\mathbb{E}[\mathbf{n}_m\mathbf{n}_n^H]$, representam a correlação entre os sinais de ruído observados por antenas pertencentes a DUs distintas. Para obter essas correlações diretamente a partir das amostras de ruído, seria necessário compartilhar informações entre as unidades, aumentando o tráfego de comunicação.

\textcolor{red}{Cada amostra $\mathbf{n}^i$ pode ser dividida em partes $\mathbf{n}_c^i$, uma para cada cluster, e a DU $c$ guarda apenas as amostras das suas próprias antenas \cite{zhao2025jsac}. Com elas, a DU $c$ calcula, sem trocar dados, só o bloco da diagonal}
\begin{equation*}
\textcolor{red}{\hat{\mathbf{R}}_{cc}=\frac{1}{N}\sum_{i=1}^{N}\mathbf{n}_c^i(\mathbf{n}_c^i)^H,}
\end{equation*}
\textcolor{red}{de tamanho $M_c \times M_c$. Os blocos fora da diagonal precisariam de amostras de outras DUs \cite{zhao2025jsac}. Neste trabalho, cada DU aplica o encolhimento de Ledoit-Wolf ao seu próprio bloco (Seção~{\hypersetup{linkcolor=red}\ref{subsec:condicionamento-R}}).}
% FONTE: paper.md:123-125 (n^i dividido em n_c^i, guardado no cluster c; R_mn = E[n_m n_n^H]);
%        Fig. 3 e paper.md:128 (blocos fora da diagonal exigem troca de dados);
%        LW por bloco: dbp_sim/m7_equalizers.py:71.

\textcolor{red}{No BCD-MMSE, a DU $c$ usa $\hat{\mathbf{R}}_{cc}$ de duas formas. Na inicialização, usa a inversa $\hat{\mathbf{R}}_{cc}^{-1}$, pois o BCD-MMSE começa pela aproximação bloco-diagonal \cite{zhao2025jsac}. Nas iterações, soma $\hat{\mathbf{R}}_{cc}$ ao termo do canal e forma a matriz $\mathbf{G}_c = E_s\mathbf{H}_c\mathbf{H}_c^H + \hat{\mathbf{R}}_{cc}$, que é invertida \cite{zhao2025jsac}. $\mathbf{G}_c$ usa $\hat{\mathbf{R}}_{cc}$, e não a sua inversa. Por isso, as duas inversões podem ser feitas em paralelo.}
% FONTE: Algoritmo 3, paper.md:415 (W^0 pela Eq. 7); Eq. 7, paper.md:136-142 (R_cc^-1);
%        Eq. 30, paper.md:504 (G_c com (1/N) sum n_c^i (n_c^i)^H).
%        G_c^-1 é calculada uma vez por PRB e reutilizada nas T iterações (não muda com l).

\textcolor{red}{Como há uma estimativa por PRB (Seção~{\hypersetup{linkcolor=red}\ref{subsec:estimativa-prb}}), cada DU calcula, a cada $T_c$, $N_{\text{PRB}}$ estimativas $\hat{\mathbf{R}}_{cc}$, $N_{\text{PRB}}$ inversões de $\hat{\mathbf{R}}_{cc}$ e $N_{\text{PRB}}$ inversões de $\mathbf{G}_c$. Com o pior caso do artigo base, $N_{\text{PRB}} = 275$ \cite{zhao2025jsac}.}
% FONTE: N_PRB = 275: paper.md:640. Contagem: planilha, tabela "Carga por PRB" (linhas 56-66).
% O que cada DU tem aceso

% É iniciado como bloco diagonal

\subsection{Aproximação bloco-diagonal (BDAC-MMSE)}

\section{Equalizador BCD-MMSE}\label{sec:bcd}
% CONTEÚDO  Roteiro T4 (formulação BCD), T5 (da eq. 26 à eq. 30) e T6 (o
%           fator a inverter, G_c = E_s H_c H_c^H + R_cc, e suas
%           propriedades: hermitiana, definida positiva).
% TEXTO     constraints/notas_redacao.md (T4-T6);
%           constraints/modelo_ruido_e_hipoteses.md, seção 10.
% ESTADO    a escrever

\subsection{Formulação por blocos}

\subsection{Atualização de cada DU}

\subsection{A matriz invertida em cada DU}\label{subsec:Gc}

\section{Inversão de matrizes}
% CONTEÚDO  Custo O(M_c^3); decomposição de Cholesky para matrizes
%           hermitianas definidas positivas; substituições direta e reversa;
%           série de Neumann e outras alternativas; por que Cholesky para G_c.
% TEXTO     Proposta do TG1, "Matrix Inversion Techniques".
% ESTADO    a escrever (traduzir e adaptar)

\subsection{Decomposição de Cholesky}

\subsection{Série de Neumann}

\subsection{Comparação}

\section{Extração de Parâmetros}
